%==============================================================================
% Research Methods - Group assignment - proposals
%==============================================================================

% This loads the article template. For a text written in English, add the
% "english" option, i.e. \documentclass[english]{hogent-article}
\documentclass{hogent-article}

% Include bibliography file
\addbibresource{references.bib}

% Info about the assignment, course and study programme

\studyprogramme{Professionele bachelor toegepaste informatica}
% TODO: replace 2N1 with the correct class group(s)!
\course{Research Methods, 2E1} 
\assignmenttype{Groepsopdracht}
\academicyear{2025-2026}
\title{Kostenoptimalisatie van AI-coding-agents in productieomgevingen: het geval van Uber}

% TODO: English text? Comment or remove the lines above and uncomment the lines
% below

% \studyprogramme{Professional bachelor in applied informatics}
% \course{Research Methods, INT-IT/2B}
% \assignmenttype{Group assignment}
% \academicyear{2025-2026}
% \title{Written preparation of research proposals}

\author{Simon Boey}
\email{simon.boey@student.hogent.be}

\author{Simon Borghart}
\email{simon.borghart@student.hogent.be}

\author{Emiel De Pelsmaeker}
\email{emiel.depelsmaeker@student.hogent.be}

\projectrepo{https://github.com/SaZoMi/RM-Group-64-Specialisatie.git}

\specialisation{AI \& Data Engineering}

\keywords{large language models, AI-agenten, kostenoptimalisatie, caching, prompt-compressie, model-routing}

\begin{document}

\begin{abstract}
  Uber besteedde in 2025 3,4 miljard dollar aan onderzoek en ontwikkeling, en toch was het volledige AI-budget van het bedrijf in april 2026 al opgebruikt door een snelle stijging in het gebruik van AI-coding-agents zoals Claude Code \autocite{Jain2026}. Dit voorbeeld illustreert een probleem dat verder reikt dan één bedrijf: bestaande technieken voor het cachen van LLM-antwoorden, zoals context caching en semantische caching, zijn ontworpen voor chatbot-achtige toepassingen en schieten tekort bij AI-agenttoepassingen. Agenttoepassingen voeren immers meerdere planningsstappen na elkaar uit en hun output hangt af van externe data en omgevingscontext, waardoor klassieke caching daar beperkt bruikbaar is \autocite{Zhang2025}. Dit zorgt voor onnodig hoge operationele kosten en latency bij bedrijven die LLM-agenten inzetten in productie. Dit onderzoeksvoorstel beschrijft een geplande studie naar welke combinatie van kostenoptimalisatietechnieken het meest geschikt is om de operationele kosten van AI-coding-agents in productie te verlagen, waarbij naast agentgeschikte caching ook prompt-compressie en model-routing onderzocht worden. Op basis van een eerste literatuurstudie wordt een probleem- en oplossingsdomein afgebakend. Daarna volgt een voorstel van methodologie waarin de drie technieken via bestaande, in de literatuur onderbouwde implementaties (Agentic Plan Caching, LLMLingua en RouteLLM) vergelijkend getest worden op een vaste set coding-taken uit SWE-bench Lite, op kosten-, latency- en kwaliteitscriteria. Het onderzoek moet resulteren in een onderbouwd handvat voor platformteams die AI-coding-agents zoals Claude Code of Cursor intern uitrollen en beheren bij grote softwareorganisaties, zoals Uber, om LLM-kosten te beheersen zonder in te boeten op de productiviteit van deze tools.
\end{abstract}

\tableofcontents

\section{Inleiding}%
\label{sec:inleiding}

In april 2026 maakte Uber bekend dat het volledige jaarbudget voor AI-tools al opgebruikt was, ondanks een totale R\&D-uitgave van 3,4 miljard dollar in 2025 \autocite{Jain2026}. Volgens Chief Technology Officer Praveen Neppalli Naga kwam dit door een sterke stijging van het gebruik van AI-coding-agents zoals Claude Code en Cursor, waarbij ongeveer 11\% van de live backend-code intussen door AI-agenten geschreven wordt \autocite{Jain2026}. Uber staat hierin niet alleen: het aantal toepassingen dat gebruikmaakt van large language models (LLM's) via commerciële API's, zoals die van OpenAI en Anthropic, of via zelf gehoste modellen, is de afgelopen jaren sterk toegenomen. Een groeiend deel van deze toepassingen zijn zogenaamde AI-agenten: systemen die zelfstandig een taak plannen, tussenstappen uitvoeren en hierbij herhaaldelijk een LLM aanroepen. Dit onderzoeksvoorstel richt zich specifiek op de operationele kostenproblematiek van dergelijke agenttoepassingen, en niet op de kwaliteit van de gegenereerde output op zich of op ethische aspecten van AI-gebruik.

Zhang et al. (2025) tonen aan dat bestaande LLM-cachingtechnieken, zoals context caching en semantische caching, primair ontworpen zijn voor het bedienen van chatbots \autocite{Zhang2025}. Bij chatbots is de output doorgaans een direct antwoord op een vraag, waardoor gelijkaardige of identieke vragen relatief eenvoudig herkend en uit een cache beantwoord kunnen worden. Bij AI-agenten, zoals de coding-agents die Uber inzet, ligt dit anders: de output van een planningsstap hangt af van externe data, zoals de bestaande codebase, en van de specifieke omgevingscontext op dat moment, waardoor twee semantisch gelijkaardige taken toch tot verschillende, niet-herbruikbare plannen kunnen leiden. Het gevolg is dat klassieke caching bij agenttoepassingen minder herbruikbaarheid oplevert dan bij chatbots, en dat de kosten en latency van agenttoepassingen daardoor onnodig hoog blijven \autocite{Zhang2025}.

Daarnaast wordt in de vakliteratuur op technieken als prompt-compressie en model-routing gewezen als aanvullende manieren om kosten te verlagen \autocite{Chen2024}. Het ontbreekt echter aan een overzicht dat deze technieken specifiek afweegt in de context van agenttoepassingen, waar de klassieke aannames van caching niet zonder meer opgaan. Daardoor hebben platformteams die AI-coding-agents zoals Claude Code op schaal uitrollen binnen een softwareorganisatie, zoals bij Uber, momenteel weinig onderbouwde houvast bij het maken van een keuze.

De specifieke doelgroep van dit onderzoek zijn platformteams en engineering-leads bij middelgrote tot grote softwareorganisaties, die verantwoordelijk zijn voor de interne uitrol en het kostenbeheer van AI-coding-agents. Meer specifiek gaat het om de teams die de infrastructuurkeuzes rond agentgebruik maken, zoals welke caching-, compressie- of routingstrategie toegepast wordt.

Dit leidt tot de volgende hoofdonderzoeksvraag:
\begin{quote}
  Welke combinatie van kostenoptimalisatietechnieken is het meest geschikt om de operationele kosten van AI-coding-agents in productie te verlagen, zonder de kwaliteit van de output significant te verminderen?
\end{quote}

Deze hoofdvraag wordt opgesplitst in deelvragen binnen het probleemdomein en het oplossingsdomein. De deelvragen zijn algemeen geformuleerd voor AI-coding-agents in een productieomgeving, en worden doorheen dit onderzoek geïllustreerd aan de hand van het geval van Uber, dat AI-coding-agents zoals Claude Code op grote schaal inzet.

Deelvragen probleemdomein:
\begin{itemize}
  \item Waarom schieten klassieke context- en semantische caching specifiek tekort bij LLM-agenttoepassingen, zoals de AI-coding-agents die bedrijven als Uber inzetten?
  \item Welke factoren bepalen de kostenstructuur van LLM-agenttoepassingen, bijvoorbeeld het aantal planningsstappen, promptlengte en modelkeuze per stap?
  \item Welke risico's ontstaan wanneer kosten worden verlaagd zonder rekening te houden met de specifieke werking van agenttoepassingen, bijvoorbeeld op vlak van kwaliteitsverlies of onbetrouwbare planning?
\end{itemize}

Deelvragen oplossingsdomein:
\begin{itemize}
  \item Hoe werkt agentgeschikte caching, zoals het cachen en hergebruiken van planstructuren in plaats van volledige antwoorden, en in welke situaties levert dit meetbare kostenbesparing op?
  \item Wat houdt prompt-compressie in en in welke mate is deze techniek toepasbaar op de planningsstappen van een agent?
  \item Hoe werkt model-routing, waarbij dynamisch een goedkoper of kleiner model wordt gekozen voor eenvoudigere planningsstappen, en welke tools of frameworks ondersteunen dit?
  \item Hoe kan de impact van deze technieken op kosten en kwaliteit gemeten en onderling vergeleken worden bij agenttoepassingen?
\end{itemize}

Het doel van dit onderzoek is een onderbouwd en vergelijkend overzicht te bieden van kostenoptimalisatietechnieken voor LLM-agenttoepassingen in productieomgevingen, gericht op platformteams bij softwareorganisaties die AI-coding-agents op schaal uitrollen, zoals Uber.

Dit onderzoeksvoorstel beperkt zich tot drie met naam benoemde technieken (agentgeschikte caching, prompt-compressie en model-routing) binnen één concreet probleemdomein, namelijk kostenbeheersing bij AI-coding-agents. Kostenbesparing, latency en kwaliteitsbehoud zijn daarbij kwantificeerbare grootheden die zowel uit de literatuur als uit de eigen proof-of-concept afgeleid kunnen worden. Zowel de literatuurstudie als de PoC zijn gebaseerd op bestaande, publiek beschikbare tools, frameworks en gepubliceerde datasets, zonder dat hiervoor toegang tot interne bedrijfsdata nodig is. Het probleem is bovendien recent aangetoond in peer-reviewed onderzoek en doet zich in de praktijk voor bij minstens één groot, publiek gedocumenteerd bedrijf. De volledige aanpak, inclusief literatuurstudie, PoC en eindredactie, is ten slotte uitgewerkt binnen de beschikbare 12 weken van dit onderzoekstraject, zoals verder toegelicht in sectie~\ref{sec:methodologie}.

De rest van dit voorstel is als volgt opgebouwd: sectie~\ref{sec:literatuurstudie} bespreekt de belangrijkste inzichten uit de literatuur, sectie~\ref{sec:methodologie} beschrijft de voorgestelde onderzoeksmethodologie, sectie~\ref{sec:resultaten} schetst de verwachte resultaten, en sectie~\ref{sec:conclusie} vat het voorstel samen.

\section{Literatuurstudie}%
\label{sec:literatuurstudie}

Zhang et al. (2025) benoemen het kernprobleem waarop dit voorstel is gebaseerd: klassieke context caching en semantische caching zijn ontwikkeld voor chatbot-toepassingen en houden geen rekening met het feit dat de output van een AI-agent afhangt van externe data en omgevingscontext \autocite{Zhang2025}. Als oplossing stellen zij Agentic Plan Caching (APC) voor, waarbij niet het volledige antwoord, maar een herbruikbare planstructuur wordt gecachet en aangepast aan de specifieke taak. Hun evaluatie op meerdere praktijktoepassingen toont een gemiddelde kostenreductie van 50,31\% en een latencyreductie van 27,28\%, bij behoud van de oorspronkelijke prestaties \autocite{Zhang2025}. Dit resultaat bevestigt zowel het bestaan van het probleem als de haalbaarheid van een oplossing die specifiek rekening houdt met de werking van agenten.

De bredere literatuur over kostenreductie bij LLM-gebruik onderscheidt daarnaast strategieën als prompt-adaptatie, modelbenadering en model-cascades \autocite{Chen2024}. Chen et al. (2024) tonen met FrugalGPT aan dat een combinatie van deze strategieën tot 98\% kostenreductie kan opleveren bij een vergelijkbare kwaliteit als het duurste beschikbare model. Deze bevindingen zijn echter voornamelijk getoetst op enkelvoudige LLM-aanroepen en niet specifiek op de meerstaps-planningsstructuur van agenttoepassingen.

Op het vlak van caching onderzoeken Regmi en Pun (2024) semantische caching. Hierbij worden niet enkel identieke, maar ook vergelijkbare vragen op basis van semantische gelijkenis uit een cache beantwoord \autocite{Regmi2024}. Liu et al. (2025) reiken hiervoor een theoretisch onderbouwd raamwerk aan voor cache-eviction onder onzekere query- en kostenverdelingen \autocite{Liu2025}. Beide werken bevestigen, net als Zhang et al. (2025), dat semantische caching in de basisvorm vooral geschikt is voor vraag-antwoordpatronen zoals bij chatbots.

Voor prompt-compressie blijft LLMLingua van Jiang et al. (2023) een veelgeciteerde referentie. Door minder relevante tokens uit de prompt te verwijderen, wordt de invoer verkleind met behoud van de essentiële betekenis \autocite{Jiang2023}. Deze techniek is in principe ook toepasbaar op de tussentijdse planningsprompts van een agent.

Op het vlak van model-routing bieden Dekoninck et al. (2024) een verenigd theoretisch kader dat routing, waarbij één model per vraag gekozen wordt, en cascading, waarbij opeenvolgend grotere modellen worden aangesproken tot een bevredigend antwoord, combineert tot een geoptimaliseerde strategie \autocite{Dekoninck2024}. Ong et al. (2024) werken dit idee concreet uit met RouteLLM, een open-source framework dat op basis van voorkeursdata leert welke vragen naar een goedkoper, kleiner model gerouteerd kunnen worden en welke een sterker model vereisen, en dat op benchmarks tot 85\% kostenreductie rapporteert bij behoud van 95\% van de kwaliteit van het sterkste model \autocite{Ong2024}. Zhen et al. (2025) plaatsen deze technieken in een breder overzicht van efficiënte LLM-inferentie, met aandacht voor zowel technieken op instantieniveau als op clusterniveau \autocite{Zhen2025}.

Uit deze literatuur blijkt dat het probleem van beperkte cachebaarheid bij agenttoepassingen recent en concreet aangetoond is, en dat prompt-compressie en model-routing potentieel aanvullende oplossingen bieden. Vergelijkend onderzoek dat deze drie technieken samen afweegt specifiek voor agenttoepassingen is echter nog beperkt beschikbaar.

Om deze drie technieken zelf op een representatieve en herhaalbare manier te kunnen vergelijken, is een gestandaardiseerde en automatisch beoordeelbare taakset nodig. Jimenez et al. (2024) ontwikkelen hiervoor SWE-bench, een benchmark van reële GitHub-issues waarbij een AI-coding-agent telkens een patch moet genereren die het issue oplost, en waarbij de correctheid van die patch automatisch getoetst wordt aan de bijhorende unit tests \autocite{Jimenez2024}. Deze benchmark is inmiddels een veelgebruikte standaard om coding-agents objectief te evalueren, en vormt de basis voor de eigen proof-of-concept in dit onderzoeksvoorstel.

\section{Methodologie}%
\label{sec:methodologie}

Om de onderzoeksvragen te beantwoorden wordt een gefaseerde aanpak voorgesteld, uitvoerbaar binnen de beschikbare 12 weken van dit onderzoekstraject (1 semester). De aanpak combineert literatuuronderzoek met een technisch, praktisch luik, zodat het voorstel niet louter theoretisch van aard blijft.

\textbf{Fase 1, probleemdomein (weken 1 tot 3).} Via literatuuronderzoek wordt geanalyseerd waarom klassieke caching tekortschiet bij agenttoepassingen en welke factoren de kostenstructuur van dergelijke toepassingen concreet bepalen. Deze fase levert antwoorden op de deelvragen van het probleemdomein en resulteert in een afgebakende probleembeschrijving met een kwantitatieve kostenanalyse van een representatief agentscenario.

\textbf{Fase 2, oplossingsdomein: literatuur (weken 4 tot 6).} Agentgeschikte caching, prompt-compressie en model-routing worden elk afzonderlijk beschreven aan de hand van de geraadpleegde vakliteratuur. Deze beschrijving wordt aangevuld met een vergelijkende analyse op basis van gepubliceerde experimentele resultaten over kostenbesparing, kwaliteitsbehoud en impact op latency, telkens getoetst aan de specifieke context van agenttoepassingen. Deze fase resulteert in een literatuurgebaseerd afwegingskader per techniek.

\textbf{Fase 3, oplossingsdomein: praktische uitwerking (weken 7 tot 10).} Om het onderzoek te onderbouwen met eigen, empirisch bewijs naast de literatuur, wordt een proof-of-concept (PoC) uitgewerkt die drie condities test op dezelfde, vaste taakset. Als taakset wordt een vaste subset van 20 opgaven uit SWE-bench Lite gebruikt, een publiek beschikbare en veelgebruikte benchmark van reële GitHub-issues waarbij een AI-coding-agent telkens een patch moet genereren die het issue oplost \autocite{Jimenez2024}. Deze taken zijn representatief voor het soort backend-codeerwerk dat coding-agents bij softwareorganisaties zoals Uber uitvoeren, en zijn dankzij de gestandaardiseerde SWE-bench-testsuite automatisch en objectief te beoordelen op correctheid (slaagt de gegenereerde patch voor de bijhorende unit tests, ja of nee).

Dezelfde 20 taken worden viermaal uitgevoerd met een bestaand, open-source agentframework (bijvoorbeeld een minimale ReAct-achtige agent-loop), telkens onder een andere conditie:
\begin{itemize}
  \item \textbf{Conditie 0, baseline:} de agent lost elke taak op zonder enige optimalisatietechniek, telkens met hetzelfde sterke model.
  \item \textbf{Conditie 1, agentgeschikte caching:} de planstructuren van de agent (de opeenvolgende tussenstappen om van issue naar patch te komen) worden gecachet en hergebruikt over gelijkaardige taken, naar het principe van Agentic Plan Caching \autocite{Zhang2025}.
  \item \textbf{Conditie 2, prompt-compressie:} de prompts die de agent bij elke planningsstap opstelt, worden voor verzending gecomprimeerd met LLMLingua, dat minder relevante tokens verwijdert met behoud van de essentiële betekenis \autocite{Jiang2023}.
  \item \textbf{Conditie 3, model-routing:} elke planningsstap van de agent wordt met RouteLLM gerouteerd naar een goedkoper, kleiner model of het oorspronkelijke sterke model, afhankelijk van de ingeschatte moeilijkheidsgraad van die stap \autocite{Ong2024}.
\end{itemize}

Per conditie wordt op dezelfde 20 taken gemeten: (1) het totale tokenverbruik (input- en outputtokens), (2) de effectieve kost in dollar op basis van de publieke prijzen van de gebruikte modellen, (3) de gemiddelde latency per opgeloste taak, en (4) de slaagkans, namelijk het percentage taken waarvoor de gegenereerde patch slaagt voor de SWE-bench-testsuite. Zo wordt per techniek een individuele kosten-batenverhouding ten opzichte van de baseline bepaald, en wordt nagegaan of de kostenbesparingen die Zhang et al. (2025), Jiang et al. (2023) en Ong et al. (2024) in hun eigen experimenten rapporteren, zich ook in deze eigen, gecontroleerde testopstelling laten reproduceren op een representatieve coding-agenttaak. Indien de tijd het toelaat, wordt in een aanvullende, vijfde conditie ook een combinatie van de drie technieken getest.

\textbf{Fase 4, synthese (weken 11 tot 12).} De bevindingen uit de literatuurstudie en de PoC worden samengebracht in een afwegingskader dat aangeeft in welke situaties welke techniek of combinatie van technieken aangewezen is voor LLM-agenttoepassingen. Deze fase omvat ook de eindredactie van het onderzoeksrapport.

Elke fase levert een concreet milestone op: een afgebakende probleembeschrijving na fase 1, een literatuurgebaseerd afwegingskader na fase 2, gemeten PoC-resultaten na fase 3, en het volledige eindrapport na fase 4. Deze opeenvolging van vier fasen van telkens twee tot vier weken past binnen de beschikbare 12 weken van dit onderzoekstraject.

\section{Verwachte resultaten}%
\label{sec:resultaten}

Het onderzoek moet een vergelijkend overzicht opleveren van kostenoptimalisatietechnieken voor LLM-agenttoepassingen. Voor elke techniek wordt een inschatting verwacht van de haalbare kostenbesparing, de invloed op kwaliteit en latency, en de situaties waarin de techniek het meest geschikt is. Daarnaast levert de proof-of-concept uit fase 3 concrete, zelf gemeten resultaten op over het tokenverbruik, de kost, de latency en de slaagkans van dezelfde 20 SWE-bench-Lite-taken onder vier condities (baseline, agentgeschikte caching, prompt-compressie en model-routing), wat toelaat de bevindingen uit de literatuur \autocite{Zhang2025, Jiang2023, Ong2024} te toetsen aan een eigen, gecontroleerde testopstelling op een representatieve coding-agenttaak. Op basis van beide bronnen wordt een afwegingskader verwacht dat platformteams, zoals dat van Uber, kan ondersteunen bij het maken van een onderbouwde keuze bij het ontwerpen van AI-coding-agents.

\section{Discussie en conclusie}%
\label{sec:conclusie}

Dit voorstel kadert een concreet en recent aangetoond probleem: klassieke LLM-cachingtechnieken schieten tekort bij AI-agenttoepassingen, waardoor de operationele kosten en latency van dergelijke toepassingen onnodig hoog blijven, zoals ook publiek geïllustreerd wordt door het budgetprobleem van Uber. Op basis hiervan wordt een onderzoek voorgesteld naar de effectiviteit van agentgeschikte caching, prompt-compressie en model-routing als oplossingsstrategieën. De voorgestelde methodologie combineert literatuuronderzoek met een eigen proof-of-concept op een vaste subset van SWE-bench Lite, waarbij telkens dezelfde 20 coding-taken onder vier condities worden getest met bestaande, in de literatuur onderbouwde implementaties (Agentic Plan Caching, LLMLingua en RouteLLM), met als doel een praktijkgericht en empirisch onderbouwd afwegingskader op te leveren voor platformteams zoals dat van Uber, binnen een tijdsbestek van 12 weken. Een mogelijke beperking van het onderzoek is de snelheid waarmee dit vakgebied evolueert, waardoor specifieke kostencijfers snel kunnen verouderen. Het voorgestelde afwegingskader is daarom opgezet om ook op langere termijn bruikbaar te blijven.

\printbibliography[heading=bibintoc]

\end{document}